

\chapter{保罗書信總結}
\section{保罗書信}


\subsection{教會書信}

%\begin{comment}
\iffalse
\def\rowi  {林前 & 所提尼 & 提摩太 &亞居拉百基拉&以弗所\\}
\def\rowii  {林後 &提摩太 & 提多及弟兄 &&馬其頓\\}
\def\rowviii {帖前 & 西拉提摩太 &&& \\}
\def\rowviiii{帖後 & 西拉提摩太 & && \\}
\def\rowiii{加拉太 & 保羅& &&\\}
\def\rowiv{以弗所 & 保羅& 推基古&&羅馬 \\}
\def\rowv{腓立比 & 提摩太& 提摩太以巴弗提&&羅馬\\}
\def\rowvi{羅馬書 & （德丟代） & 菲比&提摩太路求耶孫 &\\}
\def\rowvii{歌羅西 & 提摩太&推基古阿尼西母 &亞里達古馬可以巴弗猶士都&羅馬 \\}
\begin{table}[!hbp]
\begin{tabular}{lllll}
\toprule
書名 & 寫信人 & 送信人&同工&地點\\
\midrule
\sortfirstcol{\rowi\rowii\rowiii\rowiv\rowv\rowvi\rowvii\rowviii\rowviiii}
\bottomrule
\end{tabular}
\caption{教會書信}
\end{table}
%\end{comment}
\fi


\subsubsection{二次宣教及之前的書信}
\vspace*{-\baselineskip}
\def\rowi  {帖前5 & 西拉，提摩太&二次宣教途中&哥林多&感恩，勸勉，復活及主再來  \\}
\def\rowii {帖後3 & 西拉，提摩太 &二次宣教途中&哥林多&離道反教，沉淪之子\\}
\def\rowiii{加拉太6 & 保羅 &一次宣教之後，耶路撒冷大會之前&敘利亞的安提阿&隨從別的福音\\}
\begin{table}[!hbp]
\begin{tabular}{lllll}
\toprule
書名 & 寫信人 & 寫作時間&寫作地點&主題\\
\midrule
\sortfirstcol{\rowi\rowii\rowiii}
\bottomrule
\end{tabular}
\caption{二次宣教及之前的書信}
\end{table}

\vspace*{-\baselineskip}
\subsubsection{三次宣教途中的書信}
\vspace*{-\baselineskip}
\def\rowi  {林前16 & 所提尼 & 提摩太 &亞居拉，百基拉，司提反，福徒拿都，亞該古&以弗所&解決教會問題\\}
\def\rowii  {林後13 &提摩太 & 提多及弟兄 &&馬其頓&建立領袖同工\\}
\def\rowiii{羅馬書16 & （德丟代） & 菲比&提摩太，路求，耶孫，該猶，以拉都，括土， &哥林多&猶太人與外邦人\\}
\begin{table}[!hbp]
\begin{tabular}{llllll}
\toprule
書名 & 寫信人 & 送信人&同工&地點&主題\\
\midrule
\sortfirstcol{\rowi\rowii\rowiii}
\bottomrule
\end{tabular}
\caption{三次宣教途中的書信}
\end{table}

%\vspace*{-\baselineskip}
\subsection{監獄書——第一次在羅馬作監}
\vspace*{-\baselineskip}
\def\rowi{腓立比4 & 提摩太& 提摩太，以巴弗提&&喜樂\\}
\def\rowii{以弗所6 & 保羅& 推基古&&榮耀，教會是基督的身體 \\}
\def\rowiii{歌羅西4 & 提摩太&推基古，阿尼西母 &亞里達古，馬可，以巴弗，猶士都，路加，底馬&豐盛，基督是教會的元首 \\}
\begin{table}[!hbp]
\begin{tabular}{lllll}
\toprule
書名 & 寫信人 & 送信人&同工&主題\\
\midrule
\sortfirstcol{\rowi\rowii\rowiii}
\bottomrule
\end{tabular}
\caption{監獄書信}
\end{table}


\subsection{個人書信——第二次在羅馬作監(教牧書信)}
\begin{table}[H]
%\begin{tabularx}{12cm}{|l|l|l|l|l|X|X|}%{\textwidth}{|l|l|l|l|l|X|X|}%{12cm}{|X|X|X|X|X|X|X|}
\begin{tabularx}{17cm}{lllllll}%{\textwidth}{|l|l|l|l|l|X|X|}%{12cm}{|X|X|X|X|X|X|X|}
%\begin{tabularx}{12cm}{lllllXX}
\hline
書名 & 寫信人 &收信人&收信地點&寫作&送信人&同工\\ \hline
腓 & 提摩太 &腓利門 &歌羅西& 羅馬&推基古，阿尼西母&馬可,以巴弗,亞里達古,路加，底馬\\%\hline
提前6 & 保羅 & 提摩太 &以弗所&馬其頓&&\\%\hline
提多3 & 保羅 &提多&克里特&馬其頓&推基古，或亞提馬& \\%\hline
提後4 & 保羅 & 提摩太 &以弗所&羅馬& 推基古& 路加,友布羅、布田、利奴、革老底亞\\\hline
\end{tabularx}
\caption{個人書信}
\end{table}


\subsubsection{腓利門}
保羅盼望腓利門接納逃奴阿尼西姆。

\subsubsection{提前——解决教会领袖的问题}
\begin{itemize}%[noitemsep]
\itemsep-.25em 
\item 幾個人傳異教，講荒渺無憑的話語和無窮的家譜,虛浮的話 
\item 有人講虛浮的話,想要作教法師，卻不明白自己所講說的所論定的
\item 許米乃和亞力山大等人丟棄良心，就在真道上如同船破壞了一般 
\item 挑選監督和執事的資格 
\item 如何教導勸解教會裡不同類型的人
\item 勸提摩太如何自己做信徒的榜樣
\end{itemize}

\subsubsection{提多——解决教会会友的问题}
\begin{itemize}%[noitemsep]
\itemsep-.25em 
\item 在各城設立長老 
\item 監督和長老的資格
\item 如何在教會勸解各種類型的信徒
\end{itemize}

\subsubsection{提後——教会领袖的榜样}
\begin{itemize}%[noitemsep]
\itemsep-.25em 
\item 勸提摩太剛強壯膽，不要膽怯,不以福音為恥
\item 守住善道，做基督的精兵 
\item 末世教會光景
\end{itemize}


